\chapter*{Abstract}
\addcontentsline{toc}{chapter}{Abstract (English)}
\markboth{Abstract}{Abstract}

\label{sec:abstract-english}

This thesis compares classical and learned processing in a two-hop, single-relay communication system without a direct source--destination path. It asks where a small learned relay improves on the implemented classical baselines, and how that comparison depends on channel information and computation.

The canonical experiment uses uncoded Gray-QPSK, SISO i.i.d.\ Rayleigh fading, and receiver-side channel knowledge. Eight relay strategies are compared by Monte Carlo simulation. Learned feedforward relays improve on AF at low SNR, but DF matches or exceeds the evaluated learned relays from 6~dB upward. Replicated size sweeps identify small configurations meeting the specified degradation criteria; they do not establish an overfitting mechanism or universal minimum architecture.

Separate experiments introduce ISI, amplifier nonlinearity, and uncertain channel information. On the selected fixed three-tap BPSK channel, a 170-parameter windowed MLP substantially lowers BER relative to AF and zero-threshold DF. Matched Viterbi MLSE remains ahead over the resolved range, with target-dependent required-SNR penalties of $0.26$--$2.53$~dB. Fixed-budget MLP validation extends through 16~dB; the higher-SNR tail is not validated. In a separate composite-channel pilot sweep at 10~dB, the MLP trails pilot-aided detection at twenty pilots and leads by ten. These results concern the trained impairment families, not unseen model classes.

A separately configured coded-QPSK memory study uses a 220-parameter relay with 208 MACs per symbol; its arithmetic crosses the stated full-state MLSE count between three and five taps, while MLSE retains lower BER. Coded forwarding studies show configuration-dependent differences under the implemented decoder metrics, not a general mechanism that makes relay decoding inferior. The QPSK BCJR control measures relay-output errors only.

The evidence supports bounded benefits against specified classical pipelines. It does not establish universal learned-relay superiority, generalization outside the training families, or hardware latency and energy savings.
% \begin{center}\rule{0.5\linewidth}{0.5pt}\end{center}

% \begin{flushright}
% \begin{otherlanguage}{hebrew}
% \textbf{ארכיטקטורות למידה עמוקה לתקשורת ממסר דו-קפיצתית: מחקר השוואתי של אסטרטגיות ממסר קלאסיות ומבוססות רשתות נוירונים}
% גיל צוקרמן
% חיבור זה הוגש כמילוי חלקי של הדרישות לקבלת תואר מגיסטר למדעים (M.Sc.)
% בית הספר להנדסת חשמל — אוניברסיטת תל אביב
% 2026
% \end{otherlanguage}
% \end{flushright}

\clearpage

%% ── Acknowledgments ──────────────────────────────────────────
\cleardoublepage
\phantomsection
\chapter*{Acknowledgments}
\addcontentsline{toc}{chapter}{Acknowledgments}
\markboth{Acknowledgments}{Acknowledgments}

I would like to thank my supervisors, Dr. Anatoly Khina and Prof. Raja Giryes at Tel Aviv University, for their guidance and support throughout this research.
This work was conducted in the School of Electrical and Computer Engineering, Tel Aviv University.
